\subsection*{解方程是在寻找运动的起点}
把 $Ax=b$ 读成“哪个输入经过 $A$ 后落在 $b$”。所有可达位置是矩阵列向量的线性组合，组成列空间。若目标不在其中，怎样调输入都不会成功；目标可达之后，才轮到讨论起点是否唯一。

考虑 $(x,y)\mapsto(x+y,2x+2y)$。输出的第二分量总是第一分量的两倍，故 $(3,7)$ 不可达，而 $(3,6)$ 可达。对后者，$(3,0)$ 和 $(2,1)$ 都是起点；两输入的差 $(1,-1)$ 被送到了零。

\subsection*{零空间保存不可见的变化}
沿 $(1,-1)$ 走任意距离，一坐标增加、另一坐标减少，和不变。核描述的正是输出无法看见的输入变化。若已找到一个解 $x_0$，加任意核向量仍是解；反过来，任意另一个解减去 $x_0$，由线性性也在核里。于是整个解集是一份平移过的零空间。

逆矩阵要求每个目标恰好对应一个起点：既没有不可达目标，也没有丢失方向。在同维有限维空间中，这两方面相联系；对方阵，又与非零行列式相联系。非方阵的两种要求却可能分别成立，下一节会把它们拆开。

练习时按这个顺序思考：先判断右端是否可达，再找一个特解，最后找所有不可见的变化。这样既能解释无解和多解，也能把计算与张成、基和面积压缩接起来，而不是试图把“除以矩阵”当成通用规则。
